\begin{abstract}
  This study investigates the effects of \textit{fine-tuning} on language models tasked with navigating ASCII mazes through the use of the LLM-MRI library. By analyzing the activation graphs generated from the models' embeddings, we assess structural changes using various graph metrics with distribution inference via statistical sampling. Our findings indicate subtle, yet not significant, alterations in the graph structures after \textit{fine-tuning}, suggesting that, although \textit{fine-tuning} influences information processing, these changes are not drastically reflected in the analyzed metrics. The study highlights the importance of considering different levels of abstraction within the models to fully understand the impact of \textit{fine-tuning}. Recommended future work includes exploring diverse prompts, using \textit{embbeddings} that capture task resolution, and examining the effects of various \textit{fine-tuning} techniques and model architectures on activation graph structures.
\end{abstract}

\begin{resumo} 
  Este estudo investiga os efeitos do \textit{fine-tuning} em modelos de linguagem encarregados de navegar por labirintos em ASCII por meio da utilização da biblioteca LLM-MRI. Ao analisar os grafos de ativações gerados a partir dos embeddings dos modelos, avaliamos as mudanças estruturais utilizando diversas métricas de grafos com inferência de distribuições por amostragem estatística. Nossos achados indicam alterações sutis, porém não significativas, nas estruturas dos grafos após o \textit{fine-tuning}, sugerindo que, embora o \textit{fine-tuning} influencie o processamento de informações, essas mudanças não são drasticamente refletidas nas métricas analisadas. O estudo destaca a importância de considerar diferentes níveis de abstração dentro dos modelos para compreender plenamente o impacto do \textit{fine-tuning}. Trabalhos futuros recomendados incluem a exploração de prompts diversos, uso de \textit{embbeddings} que capturem a resolução da tarefa e os efeitos de várias técnicas de \textit{fine-tuning} e arquiteturas de modelos nas estruturas dos grafos de ativações.
\end{resumo}